\documentclass[10pt,a4paper]{amsart}
\usepackage[margin=19mm]{geometry}
\usepackage{amsmath,amssymb,mathtools}
\usepackage[scr=rsfs,cal=euler]{mathalfa}
\usepackage{xfrac}
\usepackage[unicode,hidelinks,bookmarks=false]{hyperref}
\newcommand{\ie}{i.e.\ }
\newcommand{\cf}{cf.\ }
\newcommand{\resp}{respectively\ }
\newcommand{\Subsection}{\subsection}
\providecommand{\nobreakdash}{\nobreak\hbox{-}}
\setlength{\emergencystretch}{2em}
\multlinegap0pt
\usepackage{bm}
\usepackage{bbm}
\usepackage{dsfont}
\usepackage{xspace}
\usepackage{cancel}
\usepackage{mathtools}
\usepackage{amsthm}
\makeatletter
\@ifundefined{thm}{\newtheorem{thm}{Thm}}{}
\@ifundefined{lem}{\newtheorem{lem}[thm]{Lemma}}{}
\@ifundefined{prop}{\newtheorem{prop}[thm]{Proposition}}{}
\makeatother
\makeatletter
\DeclareMathOperator{\Bl}{Bl}
\DeclareMathOperator{\Br}{Br}
\expandafter\ifx\csname MD\endcsname\relax
\newenvironment{MD}{\noindent \color{cyan}{\bf MARK:} \footnotesize}{}
\fi
\DeclareMathOperator{\Sing}{Sing}
\makeatother
\title[Brauer groups of smooth loci in linear systems and torsors o]{Brauer groups of smooth loci in linear systems and torsors over Jacobians of plane curves}
\author{Moritz Hartlieb, Weite Pi}
\date{Source: arXiv:2606.11842v2 (CC BY 4.0)}
\begin{document}
\maketitle

\noindent\emph{Source and licence.} Brauer groups of smooth loci in linear systems and torsors over Jacobians of plane curves. Version arXiv:2606.11842v2 (CC BY 4.0), distributed under the Creative Commons Attribution 4.0 International licence, as recorded in the arXiv licence metadata for this exact version and shown on the abstract page https://arxiv.org/abs/2606.11842v2. Full rights notice: licenses/arxiv-2606.11842-CC-BY-4.0.txt.\par\smallskip

\noindent\textit{Editorial scope.} Counted unit: the Lem whose source label is \texttt{lem: appendixc0} in the article (source0.tex lines 977--980, source label \texttt{lem: appendixc0}), delivered together with the article's own complete original proof of it (source0.tex lines 981--1042, one \texttt{proof} environment of 62 lines). No mathematics is written, repaired or re-derived: statement and proof are the article's own text line for line. Required context retained uncounted: prop: brauer of symm (lines 767--773). Every definition, display and label of those lines is retained unchanged. Omitted from this package and not redistributed: 1--766, 774--976, 1043--1264. Nothing inside a retained excerpt is omitted or altered: every retained body is delivered exactly as the article's lines are written, so no omission receipt of any kind accompanies this package. Citations: the retained excerpts carry no citation command at all, so no bibliography entry of the article is transcribed and no reference is redistributed. Reference adaptation: none was needed and none was made; every reference inside the delivered unit resolves inside it, so no unresolved reference or citation is produced. Printed slips kept literal: no printed word, symbol, hypothesis or number of the counted statement or of its proof was altered. Layout and class: the article's own class is \texttt{amsart} with packages including lmodern, amsmath, amsthm, amssymb, amsfonts, ulem, hyperref, mathtools, stmaryrd, biblatex, mathrsfs, graphicx, verbatim, longtable; the wrapper is the lane's pinned \texttt{amsart} editorial wrapper at 10pt on A4, which restates the theorem-like environments the retained text needs and adds the article's own macro definitions that the retained text needs (\texttt{\textbackslash Bl}, \texttt{\textbackslash Br}, \texttt{\textbackslash Sing}), with extra wrapper packages (bm, bbm, dsfont, xspace, cancel, mathtools). The delivered numbering is the unit's own. Assets: no retained span contains a figure environment, a graphics-inclusion command, a PDF-page inclusion, a table or a TikZ picture; no image file is copied into this package, the package declares no asset dependency and no image file is redistributed with it. Licence: version arXiv:2606.11842v2 of this article is distributed under the Creative Commons Attribution 4.0 International licence, as recorded in the arXiv licence metadata for this exact version and shown on the abstract page https://arxiv.org/abs/2606.11842v2. A full rights notice is delivered at licenses/arxiv-2606.11842-CC-BY-4.0.txt. Characters: every retained excerpt is pure ASCII, so the delivered engine, XeLaTeX with its default Latin Modern text font, reports no missing character.
\begin{prop}\label{prop: brauer of symm}
    For $d\geq 4$ and $n\leq d$, the Brauer group of the $n$-th symmetric product of the universal smooth plane curve of degree $d$ is isomorphic to the Brauer group of the base $|L|_{\mathrm{sm}}$, i.e.,
    $$\Br(\mathcal C_B^{(n)}) \simeq \begin{cases}0 & \text{if } d \text{ is odd,} \\
    \mathbb Z / 2 \mathbb Z &\text{if } d \text{ is even.}\end{cases}$$
More precisely, the map $\mathcal C_B^{(n)} \to B$ induces an isomorphism
$$\Br(B) \simeq \Br(\mathcal C_B^{(n)}).$$
\end{prop}

\begin{lem}\label{lem: appendixc0}
    Let $\mathcal C_B \to B \subset |\mathcal O_{\mathbb P^2}(3)|$ denote the universal smooth plane cubic. Then
    $$\Br(\mathcal C_B) \simeq \mathbb Z/2\mathbb Z.$$
\end{lem}

\begin{proof}

Note that ${\mathcal C}_{\Sigma_{\geq 2}} = \mathcal Q \cup \mathcal L$ splits into two irreducible components, since a general element $C \in \Delta_{\geq 2}$ splits as the union of a conic and line $C = Q \cup L$. We resolve the singularities of $\mathcal{C}_\Sigma$ as follows: Let $\mathbf{D} \subset \mathbb P^2 \times \mathbb P^2$ denote the diagonal, and consider the locus
$$\widetilde{\mathcal C}_{\Delta} \coloneq\overline{\{([C], p, q) \in \Delta \times ((\mathbb P^2 \times \mathbb P^2) \setminus \mathbf D) \mid p \in C,\, q \in \Sing(C)\}} \subset \Delta \times \Bl_{\mathbf{D}}(\mathbb P^2 \times \mathbb P^2).$$
The fiber of $\widetilde{\mathcal C}_\Delta \to \Bl_{\mathbf{D}}(\mathbb P^2 \times \mathbb P^2)$ over a point on the exceptional divisor, corresponding to a point $q \in \mathbb P^2$ and a tangent direction, is the locus of plane cubic curves that are singular at $q$ such that the tangent cone at $q$ contains the prescribed tangent direction. It follows from this description that the morphism
$$\widetilde{\mathcal C}_\Delta \to \Bl_{\mathbf D}(\mathbb P^2 \times \mathbb P^2)$$
is a projective bundle (of relative dimension $5$) and, in particular, $\widetilde{\mathcal C}_\Delta$ is smooth and resolves the universal singular cubic by the natural maps $\widetilde{\mathcal C}_{\Delta} \to \mathcal{C}_\Sigma \to \mathcal{C}_\Delta$, forgetting the choices of singular points.

By abuse of notation, we denote the pull-backs of $[\mathcal{Q}]$ and $[\mathcal{L}]$ to $H^2(\widetilde{\mathcal{C}}_\Delta, \mathbb Z/ 2\mathbb Z)$ in the same notations. Then, we have 
\[
H^1(\widetilde{\mathcal C}_{\Delta}, \mathbb Z / m \mathbb Z) = 0,\quad  H^2(\widetilde{\mathcal C}_\Delta, \mathbb Z / m \mathbb Z) = \langle H, h_p, h_q, E \rangle,
\]
where $H$ is the relative polarization pulled back from $|\mathcal O_{\mathbb P^2}(3)|,$ $h_p$ is the pullback of the ample generator on the copy of $\mathbb P^2$ parametrizing the (generically) non-singular point, $h_q$ is the pullback of the ample generator on the copy of $\mathbb P^2$ parametrizing the singular point, and $E$ is the pullback of the exceptional divisor on the blowup.

Since $\widetilde{\mathcal C}_{\Delta} \to \mathcal C_\Delta$ is a homeomorphism away from $\widetilde{\mathcal C}_{\Delta_{\geq 2}} \cup E$, we have
$$H^1(\mathcal C_{\Delta}, \mathbb Z / m \mathbb Z) \simeq \ker(H_{\mathrm{top}}(\widetilde{\mathcal C}_{\Delta_{\geq 2}} \cup E, \mathbb Z / m \mathbb Z) \to H_{\mathrm{top}}(\mathcal C_{\Delta_{\geq 2}} \cup \Sigma, \mathbb Z / m \mathbb Z) \oplus H^2(\widetilde{\mathcal C}_{\Delta}, \mathbb Z / m \mathbb Z)),$$
cf.\ the proof of Proposition~\ref{prop: brauer of symm},
where we recall that $\Sigma \subset \mathcal C_{\Delta}$ denotes the locus of singular points in the fibers. As before, the map $\widetilde{\mathcal{C}}_{\Delta_{\geq 2}} \cup E \to \mathcal{C}_{\Delta \geq 2} \cup \Sigma$ is generically $2$-to-$1$ on each irreducible component, so $\mathrm{Br}(\mathcal{C}_B)$ can only have $2$-primary torsion elements.

In order to compute the presentation of the classes $[\mathcal Q]$ and $[\mathcal L]$ with respect to the basis $$\langle H, h_p, h_q, E \rangle = H^2(\widetilde{\mathcal C}_{\Delta}, \mathbb Z / m \mathbb Z),$$ we consider the following pencils in $\widetilde{\mathcal C}_\Delta$:
\begin{enumerate}
\item Fix five points $p_1, \dots, p_5 \in \mathbb P^2$ and $q \in \mathbb P^2$ in general position. This datum uniquely determines a pencil of plane cubic curves that pass through the five points $p_1, \dots, p_5$ and admit a singularity at the point $q$. By marking the point $p_1$, we obtain a map
$$\varphi_H \colon \mathbb P^1 \to \widetilde{\mathcal C}_\Delta$$
satisfying
$$\deg \varphi_H^* H = 1,\;\;\; \deg \varphi_H^* h_p = \deg \varphi_H^* h_q = \deg \varphi_H^* E = 0.$$
The pencil contains $5$ reducible members that split as a conic through four of the five points and $q$ and the line through $q$ and the other point. In particular, we have
\[
\deg \varphi^*_H [\mathcal Q] = 4, \quad \deg \varphi_H^* [\mathcal L] = 1.
\]
\item Let $p_1,\ldots, p_5$ and $q$ be as above and fix a general line $l$ in $\mathbb P^2$. For any $p\in l$, there exists a unique cubic curve passing through $p_1,\ldots, p_5$ and $p$, as well as having a singularity at $q$. Marking the points $p$ and $q$, we obtain a pencil
$$\varphi_p \colon \mathbb P^1 \to \widetilde{\mathcal C}_\Delta$$
satisfying
$$\deg \varphi_p^*H = 3, \;\;\; \deg \varphi_p^*h_p = 1,  \;\;\; \deg \varphi^*_p h_q = \deg \varphi^*_p E = 0.$$
By considering the intersection between the line and the $5$ reducible members of the pencil, we obtain 
$$\deg \varphi_p^* [\mathcal Q] = 5 \cdot 2, \;\;\;\deg \varphi_p^*[\mathcal{L}] = 5.$$
\item Again, we take the same pencil, but fix a line in $\mathbb P^2$ that passes through $q$. Then, we mark the residual intersection point of this line with each member of this pencil and obtain a family
$$\varphi_E \colon \mathbb P^1 \to \widetilde{\mathcal C}_\Delta$$
satisfying
$$\deg \varphi^*_E H = \deg \varphi^*_E E = \deg \varphi^*_E h_p = 1, \;\;\; \deg \varphi_E^* h_q = 0.$$
By intersecting the $5$ reducible members of the pencil with the fixed line through $q$, we obtain
$$\deg \varphi_E^* [\mathcal Q] =  5, \;\;\; \deg \varphi^*_E [\mathcal L] = 0.$$

\item Finally, we take a different pencil as follows. Fix six points $p_1, \ldots, p_6 \in \mathbb P^2$ and a line $l$ in general position. For each $q\in l$, there exists a unique cubic curve passing through $p_1,\ldots, p_6$ and having a singularity at $q$. By marking the points $p_1$ and $q$, we obtain a map
\[
\varphi_q \colon \mathbb P^1 \to \widetilde{\mathcal{C}}_\Delta.
\]
Using classical geometry, one computes
\[
\deg \varphi_q^* H = 6, \quad \deg \varphi_q^* h_p = 0, \quad \deg \varphi_q^* h_q = 1,\quad \deg \varphi_q^* E = 0.
\]
Reducible cubics in this pencil have two types: the line $L$ passes through either one or two points among $p_1,\ldots, p_6$. From this, we obtain
\[
\deg \varphi^*_q [\mathcal Q] = 20, \quad \deg \varphi_q^* [\mathcal L] = 7.
\]

\end{enumerate}
Solving the resulting linear system of equations, it follows that 
\[
[\mathcal{Q}] = 4 H -2 h_p - 4 h_q +3 E, \quad [\mathcal{L}] = H + 2 h_p + h_q - 3E.
\]
Thus for $m=2$, the kernel we wish to compute equals $\mathbb Z/2\mathbb Z\cdot([\mathcal{Q}]-E)$, and $\mathrm{Br}(\mathcal{C}_B)[2]\simeq \mathbb Z/2\mathbb Z$. Repeating the argument with $m = 2^k$, we conclude that $\Br(\mathcal C_B) \simeq \mathbb Z/2\mathbb Z$.
\end{proof}

\end{document}
